\exo~ \textit{[Calculer]}

\begin{enumerate}
\item Dans chaque cas, déterminer un vecteur directeur de la droite d'équation cartésienne donnée.

a)~$5x+4y-2=0$~ \hspace{0.8cm }b)~$2x-9y-19=0$~ \hspace{0.8cm}c)~$y-2025=0$~~ \hspace{0.8cm}d)~$y-7x+1=0$.
 
\item Dans chaque cas, déterminer un vecteur directeur de la droite d'équation donnée et la pente si elle existe. 

a)~$3x-5y+2=0$~ \hspace{0.8cm }b)~$\dfrac{1}{2}x+\dfrac{3}{2}y-2=0$~ \hspace{0.8cm}c)~$y=6x-2$~~ \hspace{0.8cm}d)~$x=-\dfrac{1}{2}$ \hspace{0.8cm}e)~$y=\dfrac{9}{4}x-\dfrac{3}{4}$.
\end{enumerate}
  


\medskip

\exo \textit{[Représenter]}

Dans chaque cas, tracer la droite dont une équation est donnée en utilisant un vecteur directeur. Vous tracerez toutes les droites dans un même repère. 

a)~$2x-5y+2=0$~ \hspace{0.8cm }b)~$y=-2$~ \hspace{0.8cm}c)~$-x+3y+1=0$~~ \hspace{0.8cm}d)~$x=-3$ \hspace{0.8cm}e)~$y=-2x+3$.

\exo~ \textit{[Calculer, raisonner]}

\begin{enumerate}
\item \begin{enumerate}[label=\alph*)]
\item Déterminer une équation cartésienne de la droite $(d_1)$ de vecteur directeur $\vec{u}(-3;5)$ et qui passe par le point $A(5;0)$.
\item Le point de coordonnées $B(-7;21)$ appartient-il à la droite $(d_1)$?
\end{enumerate}
\item \begin{enumerate}[label=\alph*)]
\item Déterminer une équation cartésienne de la droite $(d_2)$  qui passe par les point $C(2;4)$ et $D(-1;3)$.
\item Soit E le point de coordonnées $(-7;1)$. Les points C, D et E sont-ils alignés?
\end{enumerate}
\item \begin{enumerate}[label=\alph*)]
\item Déterminer une équation cartésienne de la droite $(d_3)$  de pente $m=2$ et qui passe par le point $F(3;1)$.
\item Déterminer les coordonnées du point d'abscisse $0$ de la droite $(d_3)$.
\end{enumerate}
\end{enumerate}


\exo~ \textit{[Calculer, raisonner]}

On considère les trois droites suivantes:

 $(d_1): 3x-9y+4=0$, $(d_2): 4x-12y-1=0$ et $(d_3): 2x+3y-5=0$. 


\begin{enumerate}
\item \begin{enumerate}[label=\alph*)]
\item Étudier la position relative des droites $(d_1)$ et $(d_3)$.
\item Étudier la position relative des droites $(d_1)$ et $(d_2)$.
\item Qu'en déduisez-vous pour la position relative des droites $(d_2)$ et $(d_3$)?
\end{enumerate}
\item Calculer la pente de chacune des droites $(d_1)$, $(d_2)$ et $(d_3)$.
\item En observant les réponses des questions 1. et 2., compléter la conjecture suivante:

\og Deux droites non parallèles à l'axe des ordonnées sont parallèles si et seulement si elles ont des pentes .......... \fg{} .
\end{enumerate}

\exo~ \textit{[Calculer, raisonner]}

On considère les deux droites dont les équations cartésiennes sont les suivantes:

 $(d_1): 2x + 3y +10=0$ et $(d_2): -x+9y-5=0$. 


\begin{enumerate}
\item Montrer que $(d_1)$ et $(d_2)$ sont sécantes.
\item Justifier que les coordonnées du point d'intersection I des deux droites sont solution du système S :$\begin{cases}
2x+3y=-10 \\
-x+9y=5
\end{cases}$
\item Résoudre le système S pour déterminer les coordonnées du point I.

\end{enumerate}





